%% preamble.tex
%% This file contains the preambles for the PhD Thesis.
%% Author: Indrajit Ghosh
%% Date: Oct 21, 2024

%%----------------------------------------------------
%%		 Preamble for the project
%%-----------------------------------------------------
 
\usepackage[utf8]{inputenc}
\usepackage{graphicx}
\usepackage{enumitem} % Give extra customization on top of itemize and enumerate
\usepackage{ifthen}

\usepackage{amsmath, amsthm, amssymb} % for various mathematical operations
\usepackage{mathtools}
\usepackage{mathrsfs} % Provides \mathscr command
\usepackage{xfrac} % renders diagonal frac notation: use \sfrac{}{}
\usepackage{dsfont} % renders '1' for characteristic function...
\usepackage{ifthen}
\usepackage{emptypage}
\usepackage{parskip} % to make vertical spaces between two paragraphs
\setlength{\parindent}{15pt} % adds indentation at the starting of each chapter

\usepackage{epigraph} % Epigraphs are the pithy quotations often found at the start (or end) of a chapter
\usepackage{lipsum}

\usepackage[
    a4paper,
    bindingoffset=0.2in,
    centering,
    marginparwidth=2in,
    textwidth=6.1in,
    marginparsep=2em,
    top=2.5cm,
    bottom=2.5cm,
%    showframe
]{geometry} 

%%
%%%Header styles
%%
\usepackage{fancyhdr}
\pagestyle{fancy}
\fancyhf{}
\fancyhead[RE]{\textbf{\textit{\nouppercase{\leftmark}}}}
\fancyhead[LO]{\textbf{\textit{\nouppercase{\rightmark}}}}
\fancyfoot[C]{\thepage} 
\renewcommand{\headrulewidth}{2pt}
\renewcommand{\footrulewidth}{0pt}

%%
%%% Fancy Chapter Styles
%%
%\usepackage[Bjornstrup]{fncychap}
%Options: Conny, Bjarne, Glenn, Lenny, Sonny, Rejne, Bjornstrup

% Custom chapter heading
\makeatletter
\def\thickhrulefill{\leavevmode \leaders \hrule height 1ex \hfill \kern \z@}
\def\@makechapterhead#1{%
  \vspace*{10\p@}%
  {\parindent \z@ \raggedleft \reset@font
    \scshape \@chapapp{} \thechapter
    \par\nobreak
    \interlinepenalty\@M
    \Huge \bfseries #1\par\nobreak
    \hrulefill
    \par\nobreak
    \vskip 40\p@ % <- Reduced from 100pt to 40pt
  }}
\def\@makeschapterhead#1{%
  \vspace*{10\p@}%
  {\parindent \z@ \raggedleft \reset@font
    \scshape \vphantom{\@chapapp{} \thechapter}
    \par\nobreak
    \interlinepenalty\@M
    \Huge \bfseries #1\par\nobreak
    \hrulefill
    \par\nobreak
    \vskip 40\p@ % <- Reduced from 100pt to 40pt
  }}
\makeatother



%%%%%%%%%%%%%%%%%%%%%%%%%%%%%%%%%%%%%%%%%%%%%%%%%%%%%%%%%%%%%%%%%%%%%%%
%%%%%%%%%%%%%%%%%%%%%%%%%%%%%%%%%%%%

%%----------------------------------------------------
%%		 Theorem Styles
%%-----------------------------------------------------
%%%%
%%%% The next few commands set up the theorem type environments.
%%%% Here they are set up to be numbered section.number, but this can
%%%% be changed.
%%%%

\theoremstyle{plain}% default

\newtheorem{theorem}{Theorem}[chapter]
\newtheorem{lem}[theorem]{Lemma}
\newtheorem{prop}{Proposition}[chapter]
\newtheorem{cor}[theorem]{Corollary} % '*' produces an output without numbering

%%
%% If some other type is need, say conjectures, then it is constructed
%% by editing and uncommenting the following.
%%

%\newtheorem{conj}[theorem]{Conjecture} 


%%% 
%%% The following gives definition type environments (which only differ
%%% from theorem type environments in the choices of fonts).  The
%%% numbering is still tied to the theorem counter.
%%% 

\theoremstyle{definition}

\newtheorem{definition}{Definition}[chapter]
\newtheorem{example}{Example}[chapter]
\newtheorem{exercise}[example]{Exercise}

%%
%% Again more of these can be added by uncommenting and editing the
%% following. 
%%

\newtheorem{axiom}[definition]{Axiom}
%\newtheorem{condition}{Condition}
%\newtheorem{hypothesis}[definition]{Hypothesis}
\newtheorem{question}[definition]{Question}


%%% 
%%% The following gives remark type environments (which only differ
%%% from theorem type environments in the choices of fonts).  The
%%% numbering is still tied to the theorem counter.
%%% 

\theoremstyle{remark}

\newtheorem{remark}{Remark}[chapter]
\newtheorem*{note}{Note}
\newtheorem{case}{Case}

%% Add more here 
\newtheorem*{notation}{Notation}
%\newtheorem*{summ}{Summary}
%\newtheorem*{acknow}{Acknowledgment}

\renewcommand\qedsymbol{$\blacksquare$} % redefining the qed symbol

%%%
%%% The following, if uncommented, numbers equations within sections.
%%% 
\numberwithin{equation}{section}

%%
%%% Setting Tcolorbox for boxed 'theorems'
%%
\usepackage{tikz-cd} % Online editor: https://tikzcd.yichuanshen.de/

\usepackage[most]{tcolorbox}
% Usages: \begin{tcolorbox} Here is some text! \end{tcolorbox}

\newtcbtheorem{theo}{Theorem}{}{theorem}
%  \begin{theo}{Name} My theorem  statement \end{theo}
 

%%%%%%%%%% Bibliography management %%%%%%%%%%
\usepackage[titletoc]{appendix}

\usepackage[english]{babel}
\usepackage{csquotes}
\usepackage[nottoc]{tocbibind}
\usepackage[backend=bibtex,style=alphabetic,sorting=ynt]{biblatex}
\addbibresource{bibliography.bib}


%%% Hyper reference set up
\usepackage[bookmarksopen,bookmarksnumbered]{hyperref}
\hypersetup{
	pdftitle={\pdfTitle},
	pdfauthor={\pdfAuthor},
	pdfsubject={\pdfSubject},
	pdfcreationdate={\pdfCreationDate},
	pdfcreator={\pdfCreator},
	pdfkeywords={\pdfKeywords},
	colorlinks=\pdfColorLink,
	linkcolor={\pdfLinkColor},
	%    filecolor=magenta,      
	urlcolor=\pdfUrlColor,
	citecolor=\pdfCiteColor,
	pdfpagemode=UseOutlines,
}
\usepackage{bookmark}

\usepackage{imakeidx}
\makeindex[program=makeindex,options=-s style.ist,columns=2,intoc=true]

%%%%% Custom Commands %%%%%%%%%%%%%%
%% -----------------------------------------------------
%% Custom Macros for LaTeX
%% -----------------------------------------------------
%% Author: Indrajit Ghosh
%% Homepage: https://www.indrajitghosh.onrender.com
%% This file contains user-defined macros to enhance LaTeX functionality.
%% -----------------------------------------------------

% ########## Without math mode ##########

% Defines the Matlab symbol in small caps
\newcommand{\matlab}{{\sc Matlab} }

% Defines Holder's inequality with a properly formatted "o" character
\newcommand{\HI}{H$\ddot{\text{o}}$lder's inequality }

% Macro for bold text: usage \tb{text}
\newcommand{\tb}[1]{\textbf{#1}}

% Macro for normal text: usage \tn{text}
\newcommand{\tn}[1]{\textnormal{#1}}


% ########## Inside Math Mode ##########

% -------- Set Theory --------
% Right arrow symbol
\newcommand{\ra}{\rightarrow}

% -------- Calculus Symbols --------
% Bold vector notation
\newcommand{\cvec}[1]{{\mathbf #1}}

% Bold vector with arrow notation
\newcommand{\rvec}[1]{\vec{\mathbf #1}}

% Unit vector notation
\newcommand{\ihat}{\hat{\textbf{\i}}}
\newcommand{\jhat}{\hat{\textbf{\j}}}
\newcommand{\khat}{\hat{\textbf{k}}}

% Divergence operator
\newcommand{\mdiv}{{\rm div}}

% Gradient and curl operators
\newcommand{\grad}{\textnormal{grad}~}
\newcommand{\curl}{\textnormal{curl}~}

% -------- Linear Algebra --------
% Minor and trace functions
\newcommand{\minor}{{\rm minor}}
\newcommand{\trace}{{\rm trace}}

% Span function
\newcommand{\spn}{{\rm Span}}

% Null space, domain, and range
\newcommand{\nul}{\textnormal{\textrm{null}}}
\newcommand{\dom}{\textnormal{\textrm{dom}}}
\newcommand{\ran}{{\rm ran}}

% Homomorphism
\newcommand{\Hom}{{\rm{Hom}}}

% m x n matrices
\newcommand{\mbyn}[2]{#1\times #2}

% Tensor product and isomorphism
\newcommand{\tensor}{{\otimes}}
\newcommand{\iso}{{\cong}}

% -------- Analysis --------
% Common number sets
\newcommand{\C}{\mathbb{C}}
\newcommand{\R}{\mathbb{R}}
\newcommand{\N}{\mathbb{N}}
\newcommand{\Q}{\mathbb{Q}}
\newcommand{\Z}{\mathbb{Z}}
\newcommand{\Sp}{\mathbb{S}}
\newcommand{\F}{\mathbb{F}} % Field F

% Inner product and empty set
\newcommand{\<}{\langle}
\renewcommand{\>}{\rangle}
\renewcommand{\emptyset}{\varnothing}

% Specific vector spaces
\newcommand{\Cn}{\C^n}
\newcommand{\Zn}{\Z^n}

% -------- Operator Algebra --------
% Bounded operators and matrix algebra
\newcommand{\Bh}{\mathcal{B}(\mathcal{H})}
\newcommand{\Ch}{\mathcal{C}(\mathcal{H})}
\newcommand{\mnc}{\mathbb{M}_n(\mathbb{C})}
\newcommand{\mn}[1]{\mathbb{M}_{#1}(\mathbb{C})}

% von Neumann algebra and convergence
\newcommand{\RR}{\mathscr{R}}
\newcommand{\wkly}{\rightharpoonup}

% Norm symbols
\newcommand{\norm}[1]{\left| \left| #1 \right| \right|}
\newcommand{\opnorm}[1]{\left| \left| #1 \right| \right|_{\text{op}}}
\newcommand{\supnorm}[1]{\left| \left| #1 \right| \right|_{\text{$\infty$}}}

% Direct sum
\newcommand{\dsum}[2]{\bigoplus_{#1}^{#2}}

% Projection relationships
\newcommand{\weaker}{\precsim}
\newcommand{\strictlyweaker}{\prec}

% Murray-von Neumann equivalence
\newcommand{\mvnequiv}[1]{\stackrel{\mathclap{\normalfont\tiny\mbox{$#1$}}}{\sim}}

% -------- Mathematical Operators --------
% Distance function
\DeclareMathOperator{\dist}{dist}

% Floor and ceiling functions
\newcommand{\floor}[1]{\left\lfloor #1 \right\rfloor}
\newcommand{\ceil}[1]{\left\lceil #1 \right\rceil}

% Quotations
\newcommand{\quotes}[1]{\textquotedblleft #1\textquotedblright}

% Differential symbols
\newcommand{\diff}{\mathop{}\!d}
\newcommand{\dd}{\mathop{}\!\mathrm{d}}

% Partial derivatives
\newcommand{\pd}[2]{\frac{\partial #1 }{\partial #2}}
\newcommand{\pdS}[2]{\frac{\partial^2 #1 }{\partial #2^2}}
\newcommand{\pdT}[2]{\frac{\partial^2 #1 }{\partial #2^2}}

% Greek letters
\newcommand{\ep}{\varepsilon}

% Fraktur, calligraphy, and math script
\newcommand{\fA}{\mathfrak A}     \newcommand{\sA}{\mathcal A}		\newcommand{\kA}{\mathscr{A}}
\newcommand{\fB}{\mathfrak B}     \newcommand{\sB}{\mathcal B}		\newcommand{\kB}{\mathscr{B}}
\newcommand{\fC}{\mathfrak C}     \newcommand{\sC}{\mathcal C}		\newcommand{\kC}{\mathscr{C}}
\newcommand{\fD}{\mathfrak D}     \newcommand{\sD}{\mathcal D}		\newcommand{\kD}{\mathscr{D}}
\newcommand{\fE}{\mathfrak E}     \newcommand{\sE}{\mathcal E}		\newcommand{\kE}{\mathscr{E}}
\newcommand{\fF}{\mathfrak F}     \newcommand{\sF}{\mathcal F}		\newcommand{\kF}{\mathscr{F}}
\newcommand{\fG}{\mathfrak G}     \newcommand{\sG}{\mathcal G}		\newcommand{\kG}{\mathscr{G}}
\newcommand{\fH}{\mathfrak H}     \newcommand{\sH}{\mathcal H}		\newcommand{\kH}{\mathscr{H}}
\newcommand{\fI}{\mathfrak I}     \newcommand{\sI}{\mathcal I}		\newcommand{\kI}{\mathscr{I}}
\newcommand{\fK}{\mathfrak K}     \newcommand{\sK}{\mathcal K}		\newcommand{\kK}{\mathscr{K}}
\newcommand{\fJ}{\mathfrak J}     \newcommand{\sJ}{\mathcal J}		\newcommand{\kJ}{\mathscr{J}}
\newcommand{\fL}{\mathfrak L}     \newcommand{\sL}{\mathcal L}		\newcommand{\kL}{\mathscr{L}}
\newcommand{\fM}{\mathfrak M}     \newcommand{\sM}{\mathcal M}		\newcommand{\kM}{\mathscr{M}}
\newcommand{\fN}{\mathfrak N}     \newcommand{\sN}{\mathcal N}		\newcommand{\kN}{\mathscr{N}}
\newcommand{\fO}{\mathfrak O}     \newcommand{\sO}{\mathcal O}		\newcommand{\kO}{\mathscr{O}}
\newcommand{\fP}{\mathfrak P}     \newcommand{\sP}{\mathcal P}		\newcommand{\kP}{\mathscr{P}}
\newcommand{\fQ}{\mathfrak Q}     \newcommand{\sQ}{\mathcal Q}		\newcommand{\kQ}{\mathscr{Q}}
\newcommand{\fR}{\mathfrak R}     \newcommand{\sR}{\mathcal R}		\newcommand{\kR}{\mathscr{R}}
\newcommand{\fS}{\mathfrak S}     \newcommand{\sS}{\mathcal S}		\newcommand{\kS}{\mathscr{S}}
\newcommand{\fT}{\mathfrak T}     \newcommand{\sT}{\mathcal T}		\newcommand{\kT}{\mathscr{T}}
\newcommand{\fU}{\mathfrak U}     \newcommand{\sU}{\mathcal U}		\newcommand{\kU}{\mathscr{U}}
\newcommand{\fV}{\mathfrak V}     \newcommand{\sV}{\mathcal V}		\newcommand{\kW}{\mathscr{W}}
\newcommand{\fW}{\mathfrak W}     \newcommand{\sW}{\mathcal W}		\newcommand{\kV}{\mathscr{V}}
\newcommand{\fX}{\mathfrak X}     \newcommand{\sX}{\mathcal X}		\newcommand{\kX}{\mathscr{X}}
\newcommand{\fY}{\mathfrak Y}     \newcommand{\sY}{\mathcal Y}		\newcommand{\kY}{\mathscr{Y}}
\newcommand{\fZ}{\mathfrak Z}     \newcommand{\sZ}{\mathcal Z}		\newcommand{\kZ}{\mathscr{Z}}


%bold vector macros
\newcommand{\bu}{{\bf u}}
\newcommand{\bv}{{\bf v}}
\newcommand{\bw}{{\bf w}}
\newcommand{\bx}{{\bf x}}
\newcommand{\by}{{\bf y}}
\newcommand{\bz}{{\bf z}}
\newcommand{\bX}{{\bf X}}
\newcommand{\bY}{{\bf Y}}
\newcommand{\bZ}{{\bf Z}}
\newcommand{\btheta}{{\bf \theta}}
\newcommand{\bTheta}{{\bf \Theta}}
\newcommand{\bze}{{\bf 0}}
\newcommand{\bT}{{\bf T}}
\newcommand{\bB}{{\bf B}}
\newcommand{\bE}{{\bf E}}
\newcommand{\bN}{{\bf N}}


%%%%%%%%%%%%%%%%%%%%%%%%%%%%%%%%%%%%%%%%%%%%%%%%%%%%%%%%%%%%%%%%%%%%%%%%%%
%%%         Document Specific Commands
%%%%%%%%%%%%%%%%%%%%%%%%%%%%%%%%%%%%%%%%%%%%%%%%%%%%%%%%%%%%%%%%%%%%%%%%%%
\newcommand{\irchi}[2]{\raisebox{\depth}{$#1\chi$}} % inner command, used by \rchi
\DeclareRobustCommand{\rchi}{{\mathpalette\irchi\relax}}

\newcommand{\ipdt}[2]{\left\langle {#1, #2} \right\rangle}
\newcommand{\raR}{\textbf{R}}
\newcommand{\nulN}{\textbf{N}}
\newcommand{\graph}[1]{\text{Gr}\left( #1\right)}
\newcommand{\krein}[1]{{#1}^{\textbf{Kr}}}
\newcommand{\fried}[1]{{#1}^{\textbf{Fr}}}
\newcommand{\proj}[1]{#1_\tn{proj}} % projection lattice of an algebra
\newcommand{\sa}[1]{#1_{\tn{sa}}} % self-adjoint part of an algebra

\newcommand{\afftext}{\mathrm{aff}}
\newcommand{\aff}[1]{#1_{\afftext}}
\newcommand{\mvnaff}[1]{#1_{\afftext}^{\mathrm{MvN}}}
\newcommand{\affc}[1]{#1_{\afftext}^{c}}
\newcommand{\affs}[1]{\mathrm{Aff}_{s}\left( #1 \right)}
\newcommand{\waff}[1]{#1_{\afftext}^{\text{wor}}}
\newcommand{\Lim}{{\textnormal{\textsc{Lim}}} ~}
\newcommand{\spec}[1]{\textnormal{sp}\left( #1 \right)}

\newcommand{\cstarcat}{\mathbf{C^*Alg}}
\newcommand{\mstarcat}{\mathbf{M\sigma C^*Alg}}
\newcommand{\wstarcat}{\mathbf{W^*Alg}}

\newcommand{\mat}[2]{\mathbb{M}_#1 \left( #2 \right)}

\newcommand{\invran}{\mathrel{\overset{\makebox[0pt]{\mbox{\normalfont\tiny\sffamily inv}}}{\circ}}}

\newcommand{\blah}{{\textcolor{red}{~blah ??}}}
\newcommand{\inprogress}{{\textcolor{magenta}{\centering \LARGE \texttt{... Chapter in progress ...} \\}}}



%%%%%%%%%%% Glossaries and Acronyms %%%%%%%%%%%%%
\usepackage[%
acronym,%
symbols,%
nogroupskip,%
nonumberlist,% Uncomment for not printing the page numbers at glossaries
%sort=use
]%
{glossaries-extra}
\makenoidxglossaries
% !TEX root = ../main.tex
% =========================================================================
% symbols.tex
% Template for defining mathematical symbols for the glossary
% Author: Indrajit Ghosh
% Created On: Oct 21, 2025
% =========================================================================
%
% Usage:
%  - Define a symbol: \glsxtrnewsymbol[sort={number},description={Description}]{label}{LaTeX code}
%  - Use the symbol in text or math via: \gls{label}
%    Example: $\gls{R}$ or \gls{Bh}
%
% Notes:
%  - sort={number} determines the order in the List of Symbols
%  - description={...} is what appears in the List of Symbols
%  - label is the key used with \gls{}
%  - LaTeX code is what will be rendered in the document

% ---------------------------
% Basic Number Sets
% ---------------------------
\glsxtrnewsymbol[sort={1},description={Set of all natural numbers}]{N}{\ensuremath{\N}}
\glsxtrnewsymbol[sort={2},description={Set of all integers}]{Z}{\ensuremath{\Z}}
\glsxtrnewsymbol[sort={3},description={Set of all real numbers}]{R}{\ensuremath{\R}}
\glsxtrnewsymbol[sort={4},description={Field of complex numbers}]{C}{\ensuremath{\C}}

% ---------------------------
% Matrices and Operators
% ---------------------------
\glsxtrnewsymbol[sort={5},description={$n\times n$ matrix algebra over $\mathbb{C}$}]{mnc}{\ensuremath{\mnc}}
\glsxtrnewsymbol[sort={6},description={Algebra of bounded operators on a Hilbert space $\mathcal{H}$}]{Bh}{\ensuremath{\Bh}}
\glsxtrnewsymbol[sort={7},description={Closed unbounded operators on $\mathcal{H}$}]{Ch}{\ensuremath{\Ch}}
\glsxtrnewsymbol[sort={8},description={Graph of an operator $T$}]{grapht}{\ensuremath{\graph{T}}}

% ---------------------------
% Special Operators and Projections
% ---------------------------
\glsxtrnewsymbol[sort={9},description={Kaufman inverse of operator $T$}]{tdag}{\ensuremath{T^\dagger}}
\glsxtrnewsymbol[sort={10},description={Characteristic projection of the closed operator $T$}]{chi}{\ensuremath{\rchi(T)}}
\newcommand{\1}[1]{\mathds{1}_{#1}} % characteristic function
\glsxtrnewsymbol[sort={11},description={Indicator function of the set $\Omega$}]{indicator}{\ensuremath{\1{\Omega}}}

% ---------------------------
% von Neumann Algebra Symbols
% ---------------------------
\glsxtrnewsymbol[sort={12},description={Self-adjoint part of the von Neumann algebra $\mathscr{M}$}]{msa}{\ensuremath{\sa{\kM}}}
\glsxtrnewsymbol[sort={13},description={Lattice of projections in $\mathscr{M}$}]{projm}{\ensuremath{\proj{\kM}}}
\glsxtrnewsymbol[sort={14},description={Lattice of affiliated subspaces of $\mathscr{M}$}]{affsm}{\ensuremath{\affs{\kM}}}
\glsxtrnewsymbol[sort={15},description={Murray–von Neumann affiliated operators associated with $\mathscr{M}$}]{mvnaff}{\ensuremath{\mvnaff{\kM}}}
\glsxtrnewsymbol[sort={16},description={Densely-defined affiliated operators}]{ddaff}{\ensuremath{\aff{\kM}^{dd-c}}}
\glsxtrnewsymbol[sort={17},description={All affiliated operators associated with $\mathscr{M}$}]{maff}{\ensuremath{\aff{\kM}}}
\glsxtrnewsymbol[sort={18},description={Affiliated operators that are closed}]{maffc}{\ensuremath{\affc{\kM}}}

% ---------------------------
% C*-Algebra Symbols
% ---------------------------
\glsxtrnewsymbol[sort={19},description={Affiliated operators associated with a $C^*$-algebra $\fA$}]{aaff}{\ensuremath{\aff{\fA}}}
\glsxtrnewsymbol[sort={20},description={$n\times n$ matrices over a $C^*$-algebra $\fA$}]{mna}{\ensuremath{\mat{n}{\fA}}}

% ---------------------------
% Categories and Extensions
% ---------------------------
\glsxtrnewsymbol[sort={21},description={Category of unital $C^*$-algebras with unital *-homomorphisms}]{cstar}{\ensuremath{\cstarcat}}
\glsxtrnewsymbol[sort={22},description={Category of von Neumann algebras with normal *-homomorphisms}]{wstar}{\ensuremath{\wstarcat}}
\glsxtrnewsymbol[sort={23},description={Krein extension of the operator $S$}]{skr}{\ensuremath{\krein{S}}}
\glsxtrnewsymbol[sort={24},description={Friedrich extension of the operator $S$}]{sfr}{\ensuremath{\fried{S}}}

% ---------------------------
% Add new symbols here!
% Example:
% \glsxtrnewsymbol[sort={25},description={Your description here}]{newlabel}{\ensuremath{symbol_code}}

% =========================================================================
% Optional: include all symbols in the List of Symbols
% Remove the following line if you want only referenced symbols to appear
\glsaddall


\usepackage{acro}
% !TEX root = ../main.tex
% Abbreviations
% You can use these abbreviations in the thesis using the cmd '\ac{}'
% For example, \ac{masa}

\DeclareAcronym{masa}{
  short=MASA,
  long=Maximal Abelian Self-adjoint Algebras,
}
\DeclareAcronym{vna}{
  short=vNa,
  long=von Neumann Algebra,
}
\DeclareAcronym{mvn}{
  short=MvN,
  long=Murray-von Neumann,
}



%%%%%%%%%%%%%%%%%%%%%%%%%%%%%%%%%%%%%%%%%%%%%%%%%%%%%%




%%
%% This is the end of the preamble.
%%

    
